% !TEX root = ../main.tex

\section{Overview}

This document contains all the notes in my journey to self teach piano.

Note that the notes are words and list of concepts, ideally it won't have any detailed explanation, unless is a hard concept for me to remember. 

Sections may expand with time as I learn more concepts or details of a certain section.

\section{Scales}

\subsection{Scale for Beginners}
From the \href{https://www.youtube.com/watch?v=GugYaekR8KQ}{Video: Scale for Beginners}.

Scales Follows \textbf{Patterns:}
\begin{itemize}
    \item Major Scale: Root|W|W|H|W|W|W|H 
    \item Minor Scale: Root|W|H|W|W|H|W|W
\end{itemize}

- Patterns don't take account of the note it starts, the root note.

\textbf{Fingering} matters:

Using the right finger helps with speed, accuracy and reliability.

First Example.

\begin{definitionbox}{D-flat Mayor}
    \begin{itemize}
        \item Notes: D-flat|E-flat|F|G-flat|A-flat|B-flat|C|D-flat
        \item Fingering LH: 3|2|1|4|3|2|1|3
        \item Fingering RH: 2|3|1|2|3|4|1|2
    \end{itemize}

\end{definitionbox}

The amount of fingers that cross on each moment depend on how many notes the scale has left.

\subsection{Inversions}

